\hypertarget{appendix-f77eadcc4a93}{}\label{referencestart:appendix-f77eadcc4a93}
\section{Complex numbers and planar rotations}\label{appendix-f77eadcc4a93-complex-numbers-and-planar-rotations}

\textbf{Curriculum depth:} Introduced · \textbf{First seen:} Phase6 · \textbf{Latest development:} Phase6

\subsection{Definition and mechanism}\label{appendix-f77eadcc4a93-definition-and-mechanism}

Complex multiplication combines planar rotation and scaling. A unit complex multiplier represents a rotation; its conjugate gives the inverse. A two-component storage tuple is not automatically a rotation.

Multiplication by \(i\) sends \((x,y)\) to \((-y,x)\). Zero collapses the plane rather than representing any rotation.

\subsection{Prerequisites}\label{appendix-f77eadcc4a93-prerequisites}

\begin{itemize}
\tightlist
\item
  \hyperref[appendix-52051ad550fe]{Scalars fields and real-number models} --- Field operations define which scalar multipliers the vector laws may use.
\end{itemize}

\subsection{Used by}\label{appendix-f77eadcc4a93-used-by}

Later curriculum sources can deepen this same node; no future coverage is marked complete.

\subsection{Recall and next depth}\label{appendix-f77eadcc4a93-recall-and-next-depth}

Derive the rotation equations and separate magnitude drift from angle convention. Full production proofs, target-specific integration and learner mastery require further evidence.

\subsection{Source and navigation}\label{appendix-f77eadcc4a93-source-and-navigation}

\href{../02-Mathematics-for-Visual-Computing/main.pdf\#nameddest=01-foundations-i-ch017-rotation-first-becomes-understandable-in-a-plane}{Teaching narrative (Mathematics for Visual Computing, Mathematics for Visual Computing, Chapter 1)} · \hyperref[appendix-fd090d982a15-phase-6---parts-501-600]{Coverage and source gap} · \hyperref[appendix-a103e7846d21-phase-6---source-analysis-and-verification]{Evidence and limits} · \hyperref[appendix-4ad47bdd269a]{Dependency Map}
